\documentclass[11pt,a4paper]{article}
\usepackage[utf8]{inputenc}
\usepackage[T1]{fontenc}
\usepackage[french]{babel}
\usepackage{geometry}
\usepackage{hyperref}
\usepackage{enumitem}
\geometry{margin=2.3cm}

\title{Theorie du Tout Dissipative -- Rapport Phase I \& II \\ \large Version augmentee fractale et chiastique}
\author{Canonical Research Collective}
\date{Mars 2026}

\begin{document}
\maketitle

\section*{A. Ouverture}
Ce document augmente le manuscrit source en le replissant selon l'axe chiastique A-B-C-X-C'-B'-A'.

\section*{B. Expansion}
Les noeuds structurels extraits du manuscrit sont :
\begin{itemize}[leftmargin=1.2cm]
\item 0. Calcul initial de la marge $\mu$
\item Conjecture de Collatz -- Preuve dissipative
\item Conjecture de Goldbach -- Signature spectrale du ``Bipole Premier''
\item Conclusion
\end{itemize}

\section*{C. Matiere}
Extrait interpretable du manuscrit d'origine :
Synthese du manuscrit source : preuve dissipative et signature spectrale.

\section*{X. Centre}
Le centre chiastique est la verification numerique : convertir des calculs en STATUT de stabilite.

\section*{C'. Retour}
Le retour referme le texte sur sa fonction : preuve, seuil, verification ou acte d'unification.

\section*{B'. Miroir}
Le miroir fait correspondre architecture du manuscrit, autorite de son regime et lisibilite de sa trajectoire.

\section*{A'. Cloture}
La cloture ne remplace pas la source ; elle l'installe comme figure fractale augmentee dans l'arche canonique.

\section*{Metadonnees}
Source : dossier source interne\\
Titre detecte : \texttt{Theorie du Tout Dissipative -- Rapport Phase I \& II}\\
Auteur detecte : \texttt{Canonical Research Collective}\\
Date detectee : \texttt{Mars 2026}

\end{document}
